% -----------------------------------------------------------
\section{The Endowment}
% -----------------------------------------------------------
	\label{TEMPLE-ENDOWMENT-ENDOW}

	% ----------------------
	\subsection{Endowment-like ordinances in the Bible}
	% ----------------------
	
		

	% ----------------------
	\subsection{Endowment ceremony prior to 1990}
	% ----------------------
		\label{TEMPLE-ENDOWMENT-ENDOW-PRE-1990}
	
		(I am taking the text below from \href{http://www.ldsendowment.org}{ldsendowment.org}.)
	
		% -----
		\subsubsection{Introduction}
		% -----
			\label{TEMPLE-ENDOWMENT-ENDOW-PRE-1990-INTRO}
		
			\textit{[Initiates assemble in silence in the chapel. When all is ready, they are ushered into the Creation Room. Here and throughout the ceremony, men sit on one side of the room, women on the other.]}

			FIRST LECTURER: Brethren and sisters, we welcome you to the temple and hope you will find joy in serving in the house of the Lord this day.

			Those of you who are here to receive your own endowment should have been washed, anointed, and clothed in the garment of the holy priesthood. The ordinances of washing, anointing, and clothing in the garment of the holy priesthood, together with the ordaining on behalf of deceased brethren, were performed previously for those deceased persons whom you are representing.

			Each of you should have received a new name in connection with this company. If any of you have forgotten the new name or have not received these ordinances as explained, please stand.

			\textit{[Pause. If someone has forgotten the new name, a temple worker draws the person aside briefly to repeat the portion of the intiatory in which the new name is bestowed.]}

			Please be alert, attentive, and refrain from whispering during the presentation of the endowment. As you are asked to proceed to the veil, please do so in an orderly manner, row by row, as directed. After passing through the veil into the celestial room, and in other areas in the temple, if you need to communicate, please whisper, thus helping us maintain the quiet reverence that should prevail in the house of the Lord.

			We will now proceed with the presentation of the endowment.
			
			SECOND LECTURER: Brethren, you have been washed and pronounced clean, or that through your faithfulness you may become clean, from the blood and sins of this generation.%
				\footnote{%
					% \label{}%
					See \hyperref[DC88-74]{DC 88:74--75}; \hyperref[DC112-33]{DC 112:33} (and note that the prior verses are talking about the priesthood); 
					}
			You have been anointed to become hereafter kings and priests unto the most high God, to rule and reign in the house of Israel forever.

			Sisters, you have been washed and anointed to become queens and priestesses to your husbands.

			Brethren and sisters, if you are true and faithful, the day will come when you will be chosen, called up, and anointed kings and queens, priests and priestesses, whereas you are now anointed only to become such. The realization of these blessings depends upon your faithfulness.

			You have had a garment placed upon you, which you were informed represents the garment given to Adam when he was found naked in the garden of Eden, and which is called the garment of the holy priesthood. This you were instructed to wear throughout your life. You were informed that it will be a shield and a protection to you if you are true and faithful to your covenants.

			You have had a new name given unto you, which you were told never to divulge, nor forget. This new name is a keyword which you will be required to give at a certain place in the temple today.

			The endowment is to prepare you for exaltation in the celestial kingdom.%
				\footnote{%
					\label{ENDOWMENT-definition}%
					In 1990 Brigham Young's definition of ``endowment'' was added at this portion of the ceremony, to replace the previous sentence. The definition given comes from JD 2:31: ``Your \textit{endowment} is, to receive all those ordinances in the House of the Lord, which are necessary for you, after you have departed this life, to enable you to walk back to the presence of the Father, passing the angels who stand as sentinels, being enabled to give them the key words, the signs and tokens, pertaining to the Holy Priesthood, and gain your eternal exaltation.'' BY spoke multiple times that day, and we do have some original records from Watt (shorthand) and Thomas Bullock for other remarks he made, but I can't find any originals for this statement.
					}
			If you proceed and receive your full endowment, you will be required to take upon yourselves sacred obligations, the violation of which will bring upon you the judgments of God, for God will not be mocked.%
				\footnote{%
					% \label{}%
					See \hyperref[DC63-58]{DC 63:58} and \hyperref[DC124-71]{DC 124:71}.
					}
			
			If any of you desire to withdraw rather than accept these obligations of your own free will and choice, you may now make it known by raising your hand.

			\textit{[Pause.]}

			Brethren and sisters, as you sit here, you will hear the voices of three persons, who represent Elohim, Jehovah, and Michael. Elohim will command Jehovah and Michael to go down and organize a world. The work of the six creative periods will be represented. They will organize man in their own likeness and image, male and female. This, however, is simply figurative, so far as the man and the woman are concerned.%
				\footnote{%
					\label{ENDOWMENT-likeness}%
					See \hyperref[DC76-24]{DC 76:24}.
					}
					
		% -----
		\subsubsection{Creation}
		% -----
			\label{TEMPLE-ENDOWMENT-ENDOW-PRE-1990-CREATION}
		
			% --
			\paragraph{The first day}
			% -- 
				\label{TEMPLE-ENDOWMENT-ENDOW-PRE-1990-DAY1}
			
				\textit{[Elohim, Jehovah, and Michael are heard speaking outside the Creation Room.]}

				ELOHIM: Jehovah, Michael, see---yonder is matter unorganized. Go ye down and organize it into a world like unto the worlds that we have heretofore formed. Call your labors the first day, and bring me word.

				JEHOVAH: It shall be done, Elohim. Come, Michael, let us go down.

				MICHAEL: We will go down, Jehovah.

				JEHOVAH: Michael, see---here is matter unorganized. We will organize it into a world like unto the worlds that we have heretofore formed. We will call our labors the first day, and return and report.

				MICHAEL: We will return and report our labors of the first day, Jehovah.

				JEHOVAH: Elohim, we have been down as thou hast commanded, and have organized a world like unto the worlds that we have heretofore formed; and we have called our labors the first
				day.

				ELOHIM: It is well.
				
			% --
			\paragraph{The second day}
			% --
				\label{TEMPLE-ENDOWMENT-ENDOW-PRE-1990-DAY2}
			
				ELOHIM: Jehovah, Michael, go down again. Gather the waters together and cause the dry land to appear. The great waters call ye Seas, and the dry land call ye Earth. Form mountains and hills, great rivers and small streams, to beautify and give variety to the face of the earth. When ye have done this, call your labors the second day, and bring me word.

				JEHOVAH: It shall be done, Elohim. Come, Michael, let us go down.

				MICHAEL: We will go down, Jehovah.

				JEHOVAH: Michael, we will gather the waters together and cause the dry land to appear. The great waters we will call Seas, and the dry land we will call Earth. We will form mountains and hills, great rivers and small streams to beautify and give variety to the face of the earth. We will call our labors the second day, and return and report.

				MICHAEL: We will return and report our labors of the second day, Jehovah.

				JEHOVAH: Elohim, we have been down as thou hast commanded, and have gathered the waters together, and have caused the dry land to appear. The great waters we have called Seas, and the dry land we have called Earth. We have formed mountains and hills, great rivers and small streams, to beautify and give variety to the face of the earth; and we have called our labors the second day.

				ELOHIM: It is well.
				
			% --
			\paragraph{The third day}
			% --
				\label{TEMPLE-ENDOWMENT-ENDOW-PRE-1990-DAY3}
			
				ELOHIM: Jehovah, Michael, return again to the earth that you have organized. Divide the light from the darkness. Call the light Day, and the darkness Night. Cause the lights in the firmament to appear---the greater light to rule the day, and the lesser light to rule the night. Cause the stars also to appear and give light to the earth, the same as with other worlds heretofore created. Call your labors the third day, and bring me word.

				JEHOVAH: It shall be done, Elohim. Come, Michael, let us return again to the earth that we have organized.

				MICHAEL: We will return again, Jehovah.

				JEHOVAH: Michael, we will divide the light from the darkness, and we will call the light Day, and the darkness Night. We will cause the lights in the firmament to appear---the greater light to rule the day, and the lesser light to rule the night. We will cause the stars also to appear and give light to the earth, the same as with worlds heretofore created. We will call our labors the third day, and return and report.

				MICHAEL: We will return and report our labors of the third day, Jehovah.

				JEHOVAH: Elohim, we have been down as thou hast commanded, and have divided the light from the darkness, and have called the light Day and the darkness Night. We have caused the lights in the firmament to appear---the greater light to rule the day, and the lesser light to rule the night. We have caused the stars also to appear and give light to the earth, the same as with worlds heretofore created; and we have called our labors the third day.

				ELOHIM: It is well.
				
			% --
			\paragraph{The fourth day}
			% --
				\label{TEMPLE-ENDOWMENT-ENDOW-PRE-1990-DAY4}
			
				ELOHIM: Jehovah, Michael, return again. Place seeds of all kinds in the earth that they may spring forth as grass, flowers, shrubbery, trees, and all manner of vegetation, each bearing seed in itself after its own kind, as on the worlds we have heretofore created. Call your labors the fourth day, and bring me word.

				JEHOVAH: It shall be done, Elohim. Come, Michael, let us go down.

				MICHAEL: We will go down, Jehovah.

				JEHOVAH: Michael, we will place seeds of all kinds in the earth that they may spring forth as grass, flowers, shrubbery, trees, and all manner of vegetation, each bearing seed in itself after its own kind, as on the worlds we have heretofore created. We will call our labors the fourth day, and return and report.

				MICHAEL: We will return and report our labors of the fourth day, Jehovah.

				JEHOVAH: Elohim, we have been down as thou hast commanded, and have placed seeds of all kinds in the earth that they may spring forth as grass, flowers, shrubbery, trees, and all manner of vegetation, each bearing seed in itself, after its own kind, as on the worlds we have heretofore created; and we have called our labors the fourth day.

				ELOHIM: It is well.
				
			% --
			\paragraph{The fifth day}
			% --
				\label{TEMPLE-ENDOWMENT-ENDOW-PRE-1990-DAY5}
			
				ELOHIM: Jehovah, Michael, now that the earth is formed, divided, and beautified, and vegetation is growing thereon, return and place beasts upon the land: the elephant, the lion, the tiger, the bear, the horse, and all other kinds of animals---fowls in the air in all their varieties, fishes of all kinds in the waters, and insects and all manner of animal life upon the earth.

				Command the beasts, the fowls, the fishes, the insects, all creeping things, and other forms of animal life to multiply in their respective \hyperref[ELEMENT-1828]{elements}, each after its kind, and every kind of vegetation to multiply in its \hyperref[SPHERE-1828]{sphere}, that every form of life may fill the measure of its creation%
					\footnote{%
						\label{ENDOWMENT-MEASURE}%
						See \hyperref[DC88-19]{DC 88:19, 25}.
						} 
				and have joy therein. Call your labors the fifth day, and bring me word.

				JEHOVAH: It shall be done, Elohim. Come, Michael, let us go down.

				MICHAEL: We will go down, Jehovah.

				JEHOVAH: Michael, now that the earth is formed, divided, and beautified, and vegetation is growing thereon, we will place beasts upon the land: the elephant, the lion, the tiger, the bear, the horse, and all other kinds of animals---fowls in the air in all their varieties, fishes of all kinds in the waters, and insects and all manner of animal life upon the earth.

				We will command the beasts, the fowls, the fishes, the insects, all creeping things, and other forms of animal life to multiply in their respective elements, each after its kind, and every kind of vegetation to multiply in its sphere, that every form of life may fill the measure of its creation and have joy therein. We will call our labors the fifth day, and return and report.

				MICHAEL: It is well, Jehovah. Now that the earth is formed, divided, and beautified, with vegetation growing thereon, and provided with animal life, it is glorious and beautiful.

				JEHOVAH: It is, Michael.

				MICHAEL: Let us return and report our labors of the fifth day, Jehovah.

				JEHOVAH: Elohim, we have been down as thou hast commanded. We have placed beasts upon the land: the elephant, the lion, the tiger, the bear, the horse, and all other kinds of animals---fowls in the air in all their varieties, fishes of all kinds in the waters, and insects and all manner of animal life upon the earth.

				We have commanded the beasts, the fowls, the fishes, the insects, all creeping things, and other forms of animal life to multiply in their respective elements, each after its kind, and every kind of vegetation to multiply in its sphere, that every form of life may fill the measure of its creation and have joy therein. We have called our labors the fifth day.

				ELOHIM: It is well.
				
			% --
			\paragraph{The sixth day}
			% --
				\label{TEMPLE-ENDOWMENT-ENDOW-PRE-1990-DAY6}
				
				ELOHIM: Jehovah, Michael, is man found upon the earth?

				JEHOVAH: Man is not found on the earth, Elohim.

				ELOHIM: Jehovah, Michael, then let us go down and form man in our own likeness and in our own image, male and female, and put into him his spirit, and let us give him dominion over the beasts, the fishes, and the birds, and make him lord over the earth, and over all things on the face of the earth.

				We will plant for him a garden, eastward in Eden, and place him in it to tend and cultivate it, that he may be happy and have joy therein. We will command him to multiply and replenish the earth, that he may have joy and rejoicing in his posterity.

				We will place before him the tree of knowledge of good and evil, and we will allow Lucifer, our common enemy, whom we have thrust out, to tempt him and to try him, that he may know by his own experience the good from the evil.

				If he yields to temptation, we will give unto him the law of sacrifice, and we will provide a Savior%
					\footnote{%
						\label{ENDOWMENT-savior-first}%
						This is the first reference to Christ in the endowment, though this one isn't by name; his name appears \hyperref[TEMPLE-ENDOWMENT-ENDOW-PRE-1990-SACRIFICE]{later}, in the discussion of the law of sacrifice.
						}
				for him, as we counseled in the beginning, that man may be brought forth by the power of the redemption and the resurrection, and come again into our presence, and with us partake of eternal life and exaltation.

				We will call this the sixth day, and we will rest from our labors for a season. Come, let us go down.

				JEHOVAH: We will go down, Elohim.
				
			% --
			\paragraph{The creation of Adam and Eve}
			% --
				\label{TEMPLE-ENDOWMENT-ENDOW-PRE-1990-ADAM-EVE}
			
				\textit{[Elohim, Jehovah, and Michael enter the Creation Room. Michael sits in a chair with his head on his chest, his eyes closed, as if he were lifeless.]}

				ELOHIM: Jehovah, see the earth that we have formed. There is no man to till and take care of it. We are here to form man in our own likeness and in our own image.

				JEHOVAH: We will do so, Elohim.

				\textit{[Elohim and Jehovah pass their hands in the air over Michael's body, from head to foot.]}

				ELOHIM: Jehovah, man is now organized, and we will put into him his spirit, the breath of life, that he may become a living soul.

				\textit{[Elohim and Jehovah lift Michael's head into an upright position. He opens his eyes.]}
				
				ELOHIM: Jehovah, is it good for man to be alone?

				JEHOVAH: It is not good for man to be alone, Elohim.%
					\footnote{%
						\label{ENDOWMENT-alone}%
						In the 1931 expos\'{e} \textit{Temple Mormonism: Its Evolution, Ritual and Meaning}, Paden (the author) reports these lines as being this way (page 16):
							\begin{quote}
								Elohim---``It is not good for man to be alone.''
								
								Jehovah---``It is not, Elohim, for we are not alone.''
							\end{quote}
						If the ceremony did used to say that, then it would be a support for Adam-God, because it would mean that Jehovah was already married. That would push against the idea of Jehovah being Christ, which is one of the important parts of Adam-God.
						}

				ELOHIM: We will cause a deep sleep to come upon this man whom we have formed, and we will take from his side a rib, from which we will form a woman to be a companion and helpmeet for him.

				\textit{[Elohim and Jehovah lower Michael's head to his chest again. He closes his eyes as if asleep.]}

				JEHOVAH: Brethren and sisters, this is Michael, who helped form the earth. When he awakens from the sleep which we have caused to come upon him, he will be known as Adam and, having forgotten all, will have become as a little child. Brethren, close your eyes as if you were asleep.

				\textit{[While the brethren's eyes are closed, Elohim and Jehovah lead Eve into the Creation Room.]}

				ELOHIM: Adam, awake and arise.

				JEHOVAH: All the brethren will please arise.

				ELOHIM: Adam, here is a woman whom we have formed and whom we give unto you to be a companion and helpmeet for you. What will you call her?

				ADAM: Eve.

				ELOHIM: Why will you call her Eve?

				ADAM: Because she is the mother of all living.

				ELOHIM: That is right, Adam; because she is the mother of all living.
				
				ELOHIM: Adam, we have organized for you this earth and have planted a garden, eastward in Eden. We will place you in the garden and will there command you and Eve to multiply and replenish the earth, that you may have joy and rejoicing in your posterity.

				Jehovah, introduce Adam into the garden which we have prepared for him.

				JEHOVAH: It shall be done, Elohim.

				\textit{[Jehovah instructs the brethren to follow Adam, and the sisters to follow Eve, into the Garden Room.]}
				
		% -----
		\subsubsection{The Garden}
		% -----
			\label{TEMPLE-ENDOWMENT-ENDOW-PRE-1990-GARDEN}
		
			ELOHIM: Adam, we have created for you this earth, and have placed upon it all kinds of vegetation and animal life. We have commanded all these to multiply in their own \hyperref[SPHERE-1828]{sphere} and \hyperref[ELEMENT-1828]{element}. We give you dominion over all these things and make you, Adam, lord over the whole earth and all things on the face thereof. We now command you to multiply and replenish the earth, that you may have joy and rejoicing in your posterity.

			We have also planted for you this garden, wherein we have placed all manner of fruits, flowers, and vegetation. Of every tree of the garden thou mayest freely eat, but of the tree of knowledge of good and evil thou shalt not eat; nevertheless, thou mayest choose for thyself, for it is given unto thee. But remember that I forbid it, for in the day thou eatest thereof thou shalt surely
			die.

			Adam, remember this commandment which we have given unto you. Now go to---dress this garden, take good care of it, be happy and have joy therein.

			We shall go away, but we shall visit you again and give you further instructions.

			\textit{[Elohim and Jehovah withdraw from the Garden Room. Eve also withdraws from view, as if passing into a different portion of the garden. Adam is left alone.]}
			
			\textit{[Lucifer enters.]}

			LUCIFER: Well, Adam, you have a new world here.

			ADAM: A new world?

			LUCIFER: Yes, a new world, patterned after the old one where we used to live.

			ADAM: I know nothing about any other world.

			LUCIFER: Oh, I see--your eyes are not yet opened. You have forgotten everything. You must eat some of the fruit of this tree.

			\textit{[Lucifer pantomimes picking two pieces of fruit from the tree of knowledge of good and evil. He offers the fruit to Adam.]}

			LUCIFER: Adam, here is some of the fruit of that tree. It will make you wise.

			ADAM: I will not partake of that fruit. Father%
				\footnote{%
					\label{ENDOWMENT-father}%
					Adam (Michael) says that ``Father'' told him this; Elohim is the one who told it to him. There are various ways of understanding what Adam has said:
						\begin{compactenum}
							\item Elohim is Adam's father.
							\item Adam is referring to both Elohim and Jehovah, since in Elohim's instruction, he says ``we'' in several places.
							\item Adam doesn't mean anything special by ``father'' here, since he is currently in a state where he doesn't really know anything, and so has no other term to describe the person who gave him that instruction.
						\end{compactenum}
					Note that Eve also refers to him as ``father'' down below.
					}
			told me that in the day I should partake of it, I should surely die.

			LUCIFER: You shall not surely die, but shall be as the Gods, knowing good and evil.

			ADAM: I will not partake of it.

			LUCIFER: Oh, you will not? Well, we shall see.

			\textit{[Adam withdraws from view.]}
			
			\textit{[Eve returns.]}

			LUCIFER: Eve, here is some of the fruit of that tree. It will make you wise. It is delicious to the taste and very desirable.

			EVE: Who are you?

			LUCIFER: I am your brother.

			EVE: You, my brother, and come here to persuade me to disobey Father?

			LUCIFER: I have said nothing about Father. I want you to eat of the fruit of the tree of knowledge of good and evil that your eyes may be opened, for that is the way Father gained his knowledge. You must eat of this fruit so as to comprehend that everything has its opposite: good and evil, virtue and vice, light and darkness, health and sickness, pleasure and pain. Thus your eyes will be opened, and you will have knowledge.

			EVE: Is there no other way?

			LUCIFER: There is no other way.

			EVE: Then I will partake.

			\textit{[Eve pantomimes taking one of the pieces of fruit from Lucifer's hand and eating it.}]

			LUCIFER: There. Now go and get Adam to partake.

			\textit{[Lucifer pantomimes placing the second piece of fruit in her hand. He withdraws from view.}]

			\textit{[Adam returns.}]

			EVE: Adam, here is some of the fruit of that tree. It is delicious to the taste and very desirable.

			ADAM: Eve, do you know what fruit that is?

			EVE: Yes. It is the fruit of the tree of knowledge of good and evil.

			ADAM: I cannot partake of it. Do you not know that Father commanded us not to partake of the fruit of that tree?

			EVE: Do you intend to obey all of Father's commandments?

			ADAM: Yes, all of them.

			EVE: Do you not recollect that Father commanded us to multiply and replenish the earth? I have partaken of this fruit and by so doing shall be cast out, and you will be left a lone man in the garden of Eden.

			ADAM: Eve, I see that this must be so. I will partake that man may be.

			\textit{[Adam pantomimes eating the fruit.}]
			
			\textit{[Lucifer returns as Adam is eating the fruit.}]

			LUCIFER: That is right.

			EVE: \textit{[To Adam.}] It is better for us to pass through sorrow that we may know the good from the evil. \textit{[To Lucifer.}] I know thee now. Thou art Lucifer, he who was cast out of Father's presence for rebellion.

			LUCIFER: Yes. You are beginning to see already.

			ADAM: What is that apron you have on?

			LUCIFER: It is an emblem of my power and priesthoods.%
				\footnote{%
					\label{ENDOWMENT-apron}%
					See \hyperref[SATAN-power]{this study} on Satan's power.
					}

			ADAM: Priesthoods?

			LUCIFER: Yes, priesthoods.

			ADAM: I am looking for Father to come down to give us further instructions.

			LUCIFER: Oh, you are looking for Father to come down, are you?

			\textit{[Elohim and Jehovah are heard speaking outside the Garden Room.}]

			ELOHIM: Jehovah, we promised Adam that we would visit him and give him further instructions. Come, let us go down.

			JEHOVAH: We will go down, Elohim.

			ADAM: I hear their voices; they are coming.

			LUCIFER: See--you are naked. Take some fig leaves and make you aprons.%
				\footnote{%
					\label{ENDOWMENT-aprons}%
					\label{CHAGORAH}%
					See \hyperref[GENESIS-3-7]{Genesis 3:7}. The word ``apron'' appears only three times in the scriptures---that verse in Genesis, plus \hyperref[ACTS-19-12]{Acts 19:12} and \hyperref[MOSES-4-13]{Moses 4:13}. In Hebrew, the word for ``apron'' is \he{.ha:gOrAh}; HALOT 290b (PDF 347) says that the word has two meanings:
						\begin{compactitem}
							\item[1] \textbf{girdle}
							\item[2] \textbf{loincloth}
						\end{compactitem}
					The first definition, according to HALOT, goes with the uses in \hyperref[2SAMUEL-18-11]{2 Samuel 18:11}, \hyperref[1KINGS-2-5]{1 Kings 2:5}, \hyperref[2KINGS-3-21]{2 Kings 3:21}, and \hyperref[ISAIAH-3-24]{Isaiah 3:24}; the second meaning goes with \hyperref[GENESIS-3-7]{Genesis 3:7}. Note that this word does \textit{not} appear in connection with the priestly clothes that we read about in Exodus.
					}
			Father will see your nakedness. Quick! Hide!

			\textit{[Lucifer withdraws from view.}]

			ADAM: Brethren and sisters, put on your aprons.

			\textit{[He waits until they have done so.}]

			\textit{[To Eve.}] Come, let us hide.

			\textit{[Adam and Eve withdraw from view.}]
			
			\textit{[Elohim and Jehovah return to the Garden Room.]}

			ELOHIM: Adam!

			Adam!

			Adam, where art thou?

			\textit{[Adam returns.]}

			ADAM: I heard thy voice and hid myself, because I was naked.

			ELOHIM: Who told thee that thou wast naked? Hast thou partaken of the fruit of the tree of knowledge of good and evil, of which we commanded thee not to partake?

			ADAM: The woman thou gavest me and commanded that she should remain with me---she gave me of the fruit of the tree, and I did eat.

			ELOHIM: Eve!

			\textit{[Eve returns.]}

			What is this that thou hast done?

			EVE: The serpent beguiled me, and I did eat.
			
			ELOHIM: Lucifer!

			\textit{[Lucifer returns.}]

			ELOHIM: What hast thou been doing here?

			LUCIFER: I have been doing that which has been done on other worlds.

			ELOHIM: What is that?

			LUCIFER: I have been giving some of the fruit of the tree of knowledge of good and evil to them.

			ELOHIM: Lucifer, because thou hast done this, thou shalt be cursed above all the beasts of the field. Upon thy belly thou shalt go, and dust shalt thou eat all the days of thy life.

			LUCIFER: If thou cursest me for doing the same thing which has been done on other worlds, I will take the spirits that follow me, and they shall possess the bodies thou createst for Adam and Eve!

			ELOHIM: I will place enmity%
				\footnote{%
					% \label{}%
					Compare \hyperref[DC101-26]{DC 101:26}.
					}
			between thee and the seed of the woman. Thou mayest have power to bruise%
			\footnote{%
				% \label{}%
				Compare \hyperref[ROMANS-16-20]{Romans 16:20}, 
				}
			his heel, but he shall have power to crush thy head.

			LUCIFER: Then with that enmity I will take the treasures of the earth, and with gold and silver I will buy up armies and navies, popes and priests, and reign with blood and horror on the earth!

			ELOHIM: Depart!

			\textit{[Lucifer passes from the Garden Room into the World Room.}]
			
			ELOHIM: Jehovah, let cherubim and a flaming sword be placed to guard the way of the tree of life, lest Adam put forth his hand, and partake of the fruit thereof, and live forever in his sins.

			JEHOVAH: It shall be done, Elohim.

			\textit{[Jehovah stretches his hand towards the tree of life.}]

			Let cherubim and a flaming sword be placed to guard the way of the tree of life, lest Adam put forth his hand, and partake of the fruit thereof, and live forever in his sins.

			It is done, Elohim.
			
			% --
			\paragraph{The law of Elohim}
			% --
				\label{TEMPLE-ENDOWMENT-ENDOW-PRE-1990-ELOHIM}

				ELOHIM: %
					\footnote{%
						\label{ENDOWMENT-conception}%
						This whole paragraph was removed in the 1990 revision. There is also a short part of the next paragraph, ``because thou hast hearkened unto the voice of thy wife,'' which was removed as well.
						}%
				Eve, because thou hast hearkened to the voice of Satan, and hast partaken of the forbidden fruit, and given unto Adam, I will greatly multiply thy sorrow and thy conception. In sorrow shalt thou bring forth children; nevertheless, thou mayest be preserved in childbearing. Thy desire shall be to thy husband, and he shall rule over thee in righteousness.

				Adam, because thou hast hearkened unto the voice of thy wife and hast partaken of the forbidden fruit, the earth shall be cursed for thy sake. Instead of producing fruits and flowers spontaneously, it shall bring forth thorns, thistles, briars, and noxious weeds to afflict and torment man; and by the sweat of thy face shalt thou eat thy bread all the days of thy life, for dust thou art, and unto dust shalt thou return.

				Inasmuch as Eve was the first to eat of the forbidden fruit, if she will covenant that from this time forth she will obey your law in the Lord%
					\footnote{%
						\label{ENDOWMENT-obey-lord}%
						The 1990 revision changes this bit to ``she will obey the law of the Lord.''
						}
				and will hearken unto your counsel as you hearken unto mine, and if you will covenant that from this time forth you will obey the law of Elohim, we will give unto you the law of obedience and sacrifice, and we will provide a Savior for you, whereby you may come back into our presence and with us partake of eternal life and exaltation.

				EVE: Adam, I now covenant to obey your law as you obey our Father.%
				\footnote{%
					\label{ENDOWMENT-obey-lord-2}%
					Similar to the above change, in the 1990 revision this passage reads, ``Adam, I now covenant to obey the law of the Lord, and to hearken to your counsel as you hearken unto Father.''
					} 

				ADAM: Elohim, I now covenant with thee that from this time forth I will obey thy law and keep thy commandments.

				ELOHIM: It is well, Adam.
				
				ELOHIM: Jehovah, inasmuch as Adam and Eve have discovered their nakedness, make coats of skins as a covering for them.

				JEHOVAH: It shall be done, Elohim.

				Brethren and sisters, the garment which was placed upon you in the washing room is to cover your nakedness and represents the coat of skins spoken of. Anciently, it was made of skins. You have received the garment; also, your new name.
			
			% --
			\paragraph{The law of obedience}
			% --
				\label{TEMPLE-ENDOWMENT-ENDOW-PRE-1990-OBEDIENCE}
			
				ELOHIM: A couple will now come to the altar.

				Brethren and sisters, this couple at the altar represents all of you as if at the altar. You must consider yourselves as if you were, respectively, Adam and Eve.

				We will put the sisters under covenant to obey the law of their husbands.%
					\footnote{%
						\label{ENDOWMENT-husbands}%
						1990 revision: ``We will put each sister under covenant to obey the law of the Lord, and to hearken to the counsel of her husband, as her husband hearkens unto the counsel of the Father.''
						}
				Sisters, arise.

				Each of you bring your right arm to the square.

				\begin{quote}
					You and each of you solemnly covenant and promise before God, angels, and these witnesses at this altar that you will each observe and keep the law of your husband and abide by his counsel in righteousness.%
						\footnote{%
							\label{ENDOWMENT-husband}%
							1990 revision: ``You and each of you solemnly covenant and promise before God, angels, and these witnesses at this altar that you will each observe and keep the law of the Lord, and hearken to the counsel of your husband as he hearkens to the counsel of the Father.''
							}
				\end{quote}

				Each of you bow your head and say, ``Yes.''

				That will do.

				Brethren, arise.

				Each of you bring your right arm to the square.

				\begin{quote}
					You and each of you solemnly covenant and promise before God, angels, and these witnesses at this altar that you will obey the law of God and keep his commandments.
				\end{quote}
				
				Each of you bow your head and say, ``Yes.''

				That will do.
				
			% --
			\paragraph{The law of sacrifice}
			% --
				\label{TEMPLE-ENDOWMENT-ENDOW-PRE-1990-SACRIFICE}
			
				ELOHIM: Brethren and sisters, you are about to be put under covenant to obey and keep the law of sacrifice, as contained in the Old and New Testaments.%
					\footnote{%
						\label{ENDOWMENT-testaments}%
						The 1990 revision says ``as contained in the holy scriptures.'' A note in the source material says that the pre-1990 ceremony moves through the scriptures with the laws: ``the Old and New Testaments are named in connection with the law of sacrifice, the Book of Mormon is named in connection with the law of the gospel, and Doctrine and Covenants is named in connection with the law of consecration. The 1990 revision abandons this sequence, referring to the scriptures generically, not by specific standard works (except in the case of the law of consecration, where Doctrine and Covenants is still specifically named).''
						}
				This law of sacrifice was given to Adam in the garden of Eden, who, when he was driven out of the garden, built an altar on which he offered sacrifices.

				And after many days, an angel of the Lord appeared unto Adam, saying, ``Why dost thou offer sacrifices unto the Lord?'' And Adam said unto him, ``I know not, save the Lord commanded me.'' And then the angel spake saying, ``This thing is a similitude of the sacrifice of the Only Begotten of the Father, which is full of grace and truth. Wherefore, thou shalt do all that thou doest in the name of the Son, and thou shalt repent and call upon God in the name of the Son forevermore.''

				The posterity of Adam down to Moses, and from Moses to Jesus Christ,%
					\footnote{%
						\label{ENDOWMENT-first-christ}%
						This is the first mention of Jesus Christ by name in the endowment ceremony.
						}
				offered up the firstfruits of the field and the firstlings of the flock, which continued until the death of Jesus Christ, which ended sacrifice by the shedding of blood.

				And as Jesus Christ has laid down his life for the redemption of mankind, so we should covenant to sacrifice all that we possess, even our own lives if necessary, in sustaining and defending the kingdom of God.

				All arise.

				Each of you bring your right arm to the square.

				\begin{quote}
					You and each of you solemnly covenant and promise before God, angels, and these witnesses at this altar that you will observe and keep the law of sacrifice, as contained in the Old and New Testaments, as it has been explained to you.
				\end{quote}
				
				Each of you bow your head and say, ``Yes.''

				That will do.
				
			% --
			\paragraph{The first token of the Aaronic priesthood}
			% --
				\label{TEMPLE-ENDOWMENT-ENDOW-PRE-1990-FIRST-AARONIC}
			
				ELOHIM: We will now give unto you the first token of the Aaronic priesthood with its accompanying name, sign, and penalty.%
					\footnote{%
						\label{ENDOWMENT-penalty}%
						This is the first mention of penalties; in the 1990 revision, everything about the penalties was removed. For example, the sentence below, ``The representation of the execution of the penalties indicates different ways in which life may be taken,'' is gone, and the paragraph below this one reads ``Before doing this, we desire to impress upon your minds the sacred character of the first token of the Aaronic priesthood, with its name and sign, as well as that of all the other tokens of the holy priesthood, with their names and signs, which you will receive in the temple this day. They are most sacred and are guarded by solemn covenants and obligations, made in the presence of God, angels, and these witnesses, to hold them sacred; and under no condition will you ever divulge them, except at a certain place in the temple that will be shown you.''
						
						Here is part of a note on this from my source material: ``Originally, the penalties included graphic descriptions of the particular ways one's life would be taken in punishment for violating the covenants of non-disclosure. (The same is true of Masonic ritual.) These graphic descriptions were replaced in the 1920s with a generic statement that rather than reveal the signs, tokens, and keywords, an initiate `would suffer my life to be taken.' Note that this revision actually shifts the nature of the penalty: instead of a violent punishment falling upon those who disclose secret knowledge, the penalty becomes an affirmation that one will preserve the secret \textit{despite} threats of violence.''
						}

				Before doing this, however, we desire to impress upon your minds the sacred character of the first token of the Aaronic priesthood, with its accompanying name, sign, and penalty, as well as that of all the other tokens of the holy priesthood, with their names, signs, and penalties, which you will receive in the temple this day. They are most sacred and are guarded by solemn covenants and obligations of secrecy to the effect that under no condition, even at the peril of your life, will you ever divulge them, except at a certain place that will be shown you hereafter.

				The representation of the execution of the penalties indicates different ways in which life may be taken.

				\textit{[Elohim demonstrates the first token of the Aaronic priesthood.}]

				We give unto you the first token of the Aaronic priesthood. We desire all to receive it. All arise.

				\textit{[All initiates receive the token.}]

				The name of this token is the new name that you received in the temple today.

				\textit{[Elohim demonstrates the sign and penalty for this token.}]

				I will now explain the covenant and obligation of secrecy which are associated with this token, its name, sign, and penalty, which you will be required to take upon yourselves.

				If I were receiving my own endowment today and had been given the name of \rule{1cm}{0.15mm} as my new name, I would repeat in my mind these words after making the sign, at the same time representing the execution of the penalty:

				\begin{quote}
					I, \rule{1cm}{0.15mm}, covenant that I will never reveal the first token of the Aaronic priesthood, with its accompanying name, sign, and penalty. Rather than do so, I would suffer my life to be taken.
				\end{quote}

				\textit{[Elohim leads the initiates in repeating this oath.}]

				That will do.
			
			ELOHIM: Jehovah, see that Adam is driven out of this beautiful garden into the lone and dreary world, where he may learn from his own experience to distinguish good from evil.

			JEHOVAH: It shall be done, Elohim.

			\textit{[Jehovah instructs the brethren to follow Adam, and the sisters to follow Eve, into the World Room. Elohim and Jehovah do not accompany initiates into the World Room.}]
			
		% -----
		\subsubsection{The Telestial World}
		% -----
			\label{TEMPLE-ENDOWMENT-ENDOW-PRE-1990-TELESTIAL}
			
			ADAM: Brethren and sisters, this represents the telestial kingdom, or the world in which we now live. Adam, on finding himself in the lone and dreary world, built an altar and offered prayer, and these are the words he uttered:

			\textit{[Adam kneels at the altar and lifts his hands towards heaven, each arm raised to the square.}]

			\begin{quote}
				Oh God, hear the words of my mouth.
				
				Oh God, hear the words of my mouth.
				
				Oh God, hear the words of my mouth.%
					\footnote{%
						\label{ENDOWMENT-words-mouth}%
						The statement ``hear the words of my mouth'' does not occur in the scriptures; ``words of my mouth'' occurs ten times, nine in the OT and one in the BOM. See \hyperref[WORD-PHRASES]{here} for a complete list; four of those are most important, and they follow here:
							\begin{compactdesc}
								\item[Deuteronomy 32:1] GIVE ear, O ye heavens, and I will speak; and hear, O earth, the words of my mouth.
								\item[Psalms 19:14] Let the words of my mouth, and the meditation of my heart, be acceptable in thy sight, O LORD, my strength, and my redeemer.
								\item[Psalms 54:2] Hear my prayer, O God; give ear to the words of my mouth.
								\item[Psalms 78:1] GIVE ear, O my people, to my law: incline your ears to the words of my mouth. 
							\end{compactdesc}
						All four of these uses of the phrase are translations of the Hebrew \he{'im:rey--piy}, where the noun in construct form is \he{'emEr}, HALOT ``word'' (67a PDF 124).
						}
			\end{quote}
			
			\textit{[Lucifer comes into view.}]

			LUCIFER: I hear you. What is it you want?

			ADAM: Who are you?

			LUCIFER: I am the god of this world.%
				\footnote{%
					% \label{}%
					Compare \hyperref[2CORINTHIANS-4-4]{2 Corinthians 4:4}: ``In whom the god of this world hath blinded the minds of them which believe not, lest the light of the glorious gospel of Christ, who is the image of God, should shine unto them.'' 
					}

			ADAM: You, the god of this world?

			LUCIFER: Yes. What do you want?

			ADAM: I am looking for messengers.

			LUCIFER: Oh, you want someone to preach to you. You want religion, do you? I will have preachers here presently.%
				\footnote{%
					\label{ENDOWMENT-preachers}%
					This whole story of the preacher is removed in the 1990 revisions, and you get Lucifer saying here instead, ``There will be many willing to preach to you the philosophies of men, mingled with scripture.'' Note that a similar remark does appear down below in the pre-1990 version.
					}
			
			\textit{[A preacher enters.}]

			LUCIFER: Good morning, sir!

			PREACHER: Good morning! \textit{[Looking out over the initiates.}] A fine congregation!

			LUCIFER: Yes, they are a very good people. They are concerned about religion. Are you a preacher?

			PREACHER: I am.

			LUCIFER: Have you been to college and received training for the ministry?

			PREACHER: Certainly! A man cannot preach unless he has been trained for the ministry.

			LUCIFER: Do you preach the orthodox religion?

			PREACHER: Yes, that is what I preach.

			LUCIFER: If you will preach your orthodox religion to these people and convert them, I will pay you well.

			PREACHER: I will do my best.

			LUCIFER: \textit{[Indicating Adam.}] Here is a man who desires religion. He is very much exercised and seems to be sincere.

			PREACHER: I understand that you are inquiring after religion.

			ADAM: I was calling upon Father.

			PREACHER: I am glad to know that you were calling upon Father.

			Do you believe in a God who is without body, parts, or passions; who sits on the top of a topless throne; whose center is everywhere and whose circumference is nowhere; who fills the universe, and yet is so small that he can dwell in your heart; who is surrounded by myriads of beings who have been saved by grace, not for any act of theirs, but by his good pleasure? Do you believe in such a great being?

			ADAM: I do not. I cannot comprehend such a being.

			PREACHER: That is the beauty of it. Perhaps you do not believe in a devil, and in that great hell, the bottomless pit, where there is a lake of fire and brimstone into which the wicked are cast, and where they are continually burning but are never consumed?

			ADAM: I do not believe in any such place.

			PREACHER: My dear friend, I am sorry for you.

			LUCIFER: I am sorry, very very sorry! What is it you want?

			ADAM: I am looking for messengers from my Father.
			
			\textit{[Elohim, Jehovah, and their messengers are heard speaking outside the World Room.}]

			ELOHIM: Jehovah, send down Peter, James, and John to visit the man Adam in the telestial world, without disclosing their identity.%
				\footnote{%
					% \label{}%
					The messengers are Peter, James, and John. In \hyperref[DC7-7]{DC 7:7}, Christ speaks about them receiving some ``power and keys'' of a certain ministry until he comes. That passage is referring to their time as mortals on the earth, rather than what they experienced before that, but it's clear that they have some special authority and must have had something like that before coming to earth.
					}
			Have them observe conditions generally, see if Satan is there, and learn whether Adam has been true to the token and sign given to him in the garden of Eden. Have them then return and bring me word.

			JEHOVAH: It shall be done, Elohim.

			Peter, James, and John, go down and visit the man Adam in the telestial world, without disclosing your identity. Observe conditions generally. See if Satan is there, and learn whether Adam has been true to the token and sign given to him in the garden of Eden. Then return and bring us word.

			PETER: It shall be done, Jehovah. Come, James and John, let us go down.

			JAMES: We will go down.

			JOHN: We will go down.

			\textit{[Peter, James, and John enter the World Room.}]

			PETER: Good morning.

			LUCIFER: Good morning, gentlemen.

			PETER: What are you doing here?

			LUCIFER: Teaching religion.

			PETER: What religion do you teach?

			LUCIFER: We teach a religion made of the philosophies of men, mingled with scripture.%
				\footnote{%
					% \label{}%
					Compare the ``precepts of men'' discussed in the scriptures, such as at \hyperref[DC45-29]{DC 45:29}.
					}

			PETER: How is your religion received by this community?

			LUCIFER: Very well---excepting this man. He does not seem to believe anything we preach.

			PETER: \textit{[To Adam.}] Good morning. What do you think of the preaching of these gentlemen?

			ADAM: I cannot comprehend it.

			PETER: Can you give us some idea concerning it?

			ADAM: They preach of a God who is without body, parts, or passions; who is so large that he fills the universe, and yet is so small that he can dwell in your heart; and of a hell, without a bottom, where the wicked are continually burned but are never consumed. To me, it is a mass of confusion.

			PETER: We do not wonder that you cannot comprehend such doctrine. Have you any tokens or signs?

			LUCIFER: \textit{[Interjecting.}] Do you have any money?

			PETER: We have sufficient for our needs.

			LUCIFER: You can buy anything in this world for money.

			PETER: \textit{[To Adam.}] Do you sell your tokens or signs for money? You have them, I presume.

			ADAM: I have them, but I do not sell them for money. I hold them sacred. I am looking for the further light and knowledge Father promised to send me.

			PETER: That is right. We commend you for your integrity. Good day. We shall probably visit you again.

			\textit{[Peter, James, and John leave the World Room.}]
			
			LUCIFER: Now is the great day of my power. I reign from the rivers to the ends of the earth. There is none who dares to molest or make afraid.

			PREACHER: Shall we ever have any apostles or prophets?

			LUCIFER: No. However, there may be some who will profess revelation or apostleship. If so, just test them by asking that they perform a great miracle, such as cutting off an arm or some other member of the body and restoring it, so that the people may know that they have come with power.
			
			\textit{[Elohim, Jehovah, and their messengers are heard speaking outside the World Room.}]

			PETER: Jehovah, we have visited the man Adam in the telestial world as thou didst command us. We found Satan there, with his ministers, preaching all manner of false doctrine and striving to lead the posterity of Adam astray. But Adam has been true and faithful to the token and sign given him in the garden of Eden and is waiting for the further light and knowledge you promised to send him. This is our report.

			JEHOVAH: It is well, Peter, James, and John.

			Elohim---Peter, James, and John have been down to the man Adam in the telestial world. They found Satan there, with his ministers, preaching all manner of false doctrine and striving to lead the posterity of Adam astray. But Adam has been true and faithful to the token and sign given him in the garden of Eden, and he is waiting for the further light and knowledge you promised to send him. This is their report.

			ELOHIM: It is well.

			Jehovah, instruct Peter, James, and John to go down in their true character, as apostles of the Lord Jesus Christ, to the man Adam and his posterity in the telestial world and to cast Satan out of their midst. Instruct them to give unto Adam and his posterity the law of the gospel as contained in the Book of Mormon and the Bible; also, a charge to avoid all lightmindedness, loud laughter, evil speaking of the Lord's anointed, the taking of the name of God in vain, and every other unholy and impure practice; and cause these to be received by covenant.

			Instruct Peter, James, and John further to clothe Adam and his posterity in the robes of the holy priesthood, with the robe on the left shoulder, and to give unto them the second token of the Aaronic priesthood, with its accompanying name, sign, and penalty. Then have them return and bring me word.

			JEHOVAH: It shall be done, Elohim.

			Peter, James, and John, go down in your true character, as apostles of the Lord Jesus Christ,%
				\footnote{%
					\label{ENDOWMENT-character}%
					As I mention \hyperref[ADAMGOD-TEMPLE]{here}, it's weird but not completely absurd for Jehovah to speak to the apostles this way if Jehovah is Christ.
					}
			to the man Adam and his posterity in the telestial world. Cast Satan out of their midst. Give unto them the law of the gospel as contained in the Book of Mormon and the Bible; also, a charge to avoid all lightmindedness, loud laughter, evil speaking of the Lord's anointed, the taking of the name of God in vain, and every other unholy and impure practice. Cause them to receive these by covenant.

			Clothe them in the robes of the holy priesthood, with the robe on the left shoulder, and give unto them the second token of the Aaronic priesthood, with its accompanying name, sign, and penalty. Then return and bring us word.

			PETER: It shall be done, Jehovah. Come, James and John, let us go down.
			
			\textit{[Peter, James, and John return to the World Room.}]

			PETER: I am Peter.

			JAMES: I am James.

			JOHN: I am John.

			LUCIFER: Yes, I thought I knew you. \textit{[To the preacher.}] Do you know who these men are? They claim to be apostles. Try them!

			PREACHER: \textit{[To Peter.}] Do you profess to be an apostle of the Lord Jesus Christ?

			PETER: We do.

			PREACHER: This man told me that we should never have any revelation or apostles, but if any should come professing to be apostles, I was to ask them to cut off an arm or some other member of the body and then restore it, so that the people might know that they came with power.

			PETER: We do not satisfy men's curiosity in that manner. It is a wicked and an adulterous generation that seeks for a sign. Do you know who that man is? He is Satan.

			PREACHER: What? The devil?

			PETER: That is one of his names.

			PREACHER: He is quite a different person from what he told me the devil is. He said the devil has claws like a bear's on his hands, horns on his head, and a cloven foot, and that when he speaks he has the roar of a lion.

			PETER: He has said this to deceive you, and I would advise you to get out of his employ.

			PREACHER: Your advice is good. But if I leave his employ, what will become of me?

			PETER: We will preach the gospel unto you, with the rest of Adam's posterity.

			PREACHER: That is good. \textit{[To Lucifer.}] I would like to have a settlement. I want you to pay me for preaching.

			LUCIFER: I am ready to keep my word and fulfill my part of the agreement. I promised to pay you if you would convert these people, and they have nearly converted you! You can get out of my kingdom---I want no such men in it!

			\textit{[The preacher exits.}]
			
			LUCIFER: \textit{[To Peter.}] Now what are you going to do?

			PETER: We will dismiss you without further argument.

			LUCIFER: Ah! You have looked over my kingdom and my greatness and glory.%
				\footnote{%
					% \label{}%
					Compare \hyperref[LUKE-4-5]{Luke 4:5--8}, especially verse 6, where the tempter refers to \gr{ἐξουσία}.
					} 
			Now you want to take possession of the whole of it.

			\textit{[He indicates the initiates.}]

			I have a word to say concerning these people. If they do not walk up to every covenant they make at these altars in this temple this day, they will be in my power!

			PETER: Satan, we command you to depart.

			LUCIFER: By what authority?

			PETER: \textit{[Raising his right arm to the square.}] In the name of Jesus Christ, our Master.

			\textit{[Lucifer exits.}]
			
			PETER: Adam, we are true messengers from the Father and have come down to give unto you the further light and knowledge he promised to send you.

			ADAM: How shall I know that you are true messengers?

			PETER: By our giving unto you the token and sign you received in the garden of Eden.

			\textit{[Peter takes Adam by the hand in the first token of the Aaronic priesthood.}]

			ADAM: What is that?

			PETER: The first token of the Aaronic priesthood.

			ADAM: Has it a name?

			PETER: It has.

			ADAM: Will you give it to me?

			PETER: I cannot, for it is the new name. But this is the sign, and this represents the execution of the penalty.

			\textit{[Peter demonstrates the sign and penalty as he says this.}]

			ADAM: Now I know that you are true messengers sent down from Father.

			\textit{[To the initiates.}] These are true messengers. I exhort you to give strict heed to their counsel and teachings, and they will lead you in the way of life and salvation.
			
			% ---
			\paragraph{The law of the gospel}
			% ---
				\label{TEMPLE-ENDOWMENT-ENDOW-PRE-1990-GOSPEL}
				
				PETER: A couple will now come to the altar.%
					\footnote{%
						\label{ENDOWMENT-peter-altar}%
						Notice that Peter has replaced Elohim at the altar; we aren't interacting directly with the Father any longer.
						}

				Brethren and sisters, this couple at the altar represent all of you as if at the altar, and you will be under the same obligations as they will be.

				We are required to give unto you the law of the gospel%
					\footnote{%
						% \label{}%
						See \hyperref[DC88-78]{DC 88:78}.
						}
				as contained in the Book of Mormon and the Bible;%
					\footnote{%
						\label{ENDOWMENT-law-gospel}%
						It's a little weird that the endowment doesn't explain the law of the gospel more. What exactly is it? What are we supposed to understand about it?
						}
				to give unto you, also, a charge to avoid all lightmindedness,%
					\footnote{%
						% \label{}%
						See \hyperref[DC84-54]{DC 84:54} and \hyperref[DC84-61]{61}, \hyperref[DC88-69]{DC 88:69}, \hyperref[DC88-121]{DC 88:121} (in this verse, it's clear that this instruction only applies to a limited number of people in restricted circumstances, I think),
						}
				loud laughter,%
					\footnote{%
						% \label{}%
						See \hyperref[DC88-69]{DC 88:69}, where the Lord instructs those listening to ``cast away...your excess of laughter.'' See also \hyperref[1NEPHI-18-9]{1 Nephi 18:9}, \hyperref[DC88-121]{DC 88:121},
						}
				evil speaking of the Lord's anointed, the taking of the name of God in vain, and every other unholy and impure practice; and to cause you to receive these by covenant.

				All arise.

				Each of you bring your right arm to the square.

				\begin{quote}
					You and each of you covenant and promise before God, angels, and these witnesses at this altar that you will observe and keep the law of the gospel and this charge, as it has been explained to you.
				\end{quote}
				
				Each of you bow your head and say, ``Yes.''

				That will do.
				
			% ---
			\paragraph{Robes of the Aaronic priesthood}
			% ---
				\label{TEMPLE-ENDOWMENT-ENDOW-PRE-1990-ROBES-AARONIC}

				PETER: We are instructed to clothe you in the robes of the holy priesthood.

				Place the robe on your left shoulder. Place the cap on your head with the bow over the right ear. Replace the apron. Tie the girdle with the bow on the right side. Remove the slippers from your feet, and put them on again as part of the temple clothing.

				You may now proceed to clothe.
				
			% ---
			\paragraph{The second token of the Aaronic priesthood}
			% ---
				\label{TEMPLE-ENDOWMENT-ENDOW-PRE-1990-SECOND-AARONIC}
				
				PETER: A couple will now come to the altar.

				With the robe on the left shoulder, you are prepared to officiate in the ordinances of the Aaronic priesthood.%
					\footnote{%
						\label{ENDOWMENT-officiate-aaronic}%
						See \hyperref[DC107-14]{DC 107:14, 20}, and discussion of ``outward ordinances'' \hyperref[ORDINANCE-PHRASES]{here}. Peter says that both the men and women are prepared to ``officiate in the ordinances of the Aaronic priesthood''---usually we would think of those ordinances as involving baptism and the sacrament. He could be referring to those, and could also be referring to temple-specific ordinances governed by the Aaronic priesthood.
						}

				We will now give unto you the second token of the Aaronic priesthood with its accompanying name, sign, and penalty.

				\textit{[Peter demonstrates the token.}]

				We desire all to receive it. All arise.

				\textit{[All initiates receive the token. Peter then reveals the name, sign, and penalty for this token. Not all initiates receive the same name.}]

				I will now explain the covenant and obligation of secrecy which are associated with this token, its name, sign, and penalty, which you will be required to take upon yourselves.

				If I were receiving my own endowment today, and if \textit{[here Peter provides a sample name for the token}], I would repeat in my mind these words, after making the sign, at the same time representing the execution of the penalty:

				\begin{quote}
					I, \rule{1cm}{0.15mm}, covenant that I will never reveal the second token of the Aaronic priesthood, with its accompanying name, sign, and penalty. Rather than do so, I would suffer my life to be taken.
				\end{quote}
				
				\textit{[Peter leads the initiates in repeating this oath.}]

				That will do.

				We will return and report.

				\textit{[Peter, James, and John exit.}]
				
				\textit{[Elohim, Jehovah, and their messengers are heard speaking outside the World Room.}]

				PETER: Jehovah, we have been down to the man Adam and his posterity in the telestial world and have cast Satan out of their midst. We have given unto them the law of the gospel as contained in the Book of Mormon and the Bible; also, a charge to avoid all lightmindedness, loud laughter, evil speaking of the Lord's anointed, the taking of the name of God in vain, and every other unholy and impure practice, and have caused them to receive these by covenant.

				We have also clothed them in the robes of the holy priesthood and have given unto them the second token of the Aaronic priesthood, with its accompanying name, sign, and penalty. This is our report.

				JEHOVAH: It is well, Peter, James, and John.

				Elohim---Peter, James, and John have been down to the man Adam and his posterity in the telestial world, have cast Satan out of their midst, and have done all else that they were commanded to do.

				ELOHIM: It is well.

				Jehovah, send down Peter, James, and John again to the telestial world. Have Adam and his posterity change their robes to the right shoulder, preparatory to officiating in the ordinances of the Melchizedek priesthood, and introduce them into the terrestrial world.

				Instruct Peter, James, and John further to give unto Adam and his posterity the law of chastity and to put them under covenant to obey this law, which is that the daughters of Eve and the sons of Adam shall have no sexual intercourse except with their husbands or wives to whom they are legally and lawfully wedded, and to give unto them the first token of the Melchizedek priesthood, or Sign of the Nail, with its accompanying name, sign, and penalty. Have them return and bring me word.

				JEHOVAH: It shall be done, Elohim.

				Peter, James and John, go down again to the telestial world. Instruct Adam and his posterity to change their robes to the right shoulder, preparatory to officiating in the ordinances of the Melchizedek priesthood, and introduce them into the terrestrial world.

				Give unto Adam and his posterity the law of chastity and put them under covenant to obey this law, which is that the daughters of Eve and the sons of Adam shall have no sexual intercourse except with their husbands or wives to whom they are legally and lawfully wedded. Give unto them the first token of the Melchizedek priesthood, or Sign of the Nail, with its accompanying name, sign, and penalty; and return and bring us word.

				PETER: It shall be done, Jehovah. Come, James and John, let us go down.
				
			% ---
			\paragraph{Robes of the Melchizedek priesthood}
			% ---
				\label{TEMPLE-ENDOWMENT-ENDOW-PRE-1990-ROBES-MELCHIZEDEK}
				
				\textit{[Peter, James, and John return to the World Room.}]

				PETER: We are instructed to have you change the robe to the right shoulder, preparatory to officiating in the ordinances of the Melchizedek priesthood, and to introduce you into the terrestrial world.

				You may now make the change by removing the robe.

				\textit{[The initiates change the robe to the right shoulder. They then pass into the Terrestrial Room.}]
		% -----
		\subsubsection{The Terrestrial World}
		% -----
			\label{TEMPLE-ENDOWMENT-ENDOW-PRE-1990-TERRESTRIAL}
			
			% --
			\paragraph{The law of chastity}
			% --
				\label{TEMPLE-ENDOWMENT-ENDOW-PRE-1990-CHASTITY}
				
				PETER: A couple will now come to the altar.

				We are instructed to give unto you the law of chastity. This I will explain. To the sisters, it is that no one of you will have sexual intercourse except with your husband to whom you are legally and lawfully wedded. To the brethren it is that no one of you will have sexual intercourse except with your wife to whom you are legally and lawfully wedded.%
					\footnote{%
						\label{ENDOWMENT-lawfully}%
						In the 1990 revision, men and women receive this covenant at the same time, and they agree that they will not have sexual relations (not ``intercourse'') except with their ``husband or wife to whom they are legally and lawfully wedded.'' Note that this way of presenting the covenant makes it possible for a person in a gay relationship to obey the covenant, because it isn't specific enough to exclude that relationship with its current wording (I'm writing this in 2021; I haven't been to the endowment for a while and so it could be different now.)
						}

				Sisters, please arise.

				Each of you bring your right arm to the square.

				\begin{quote}
					You and each of you covenant and promise before God, angels, and these witnesses at this altar that you will observe and keep the law of chastity, as it has been explained to you.
				\end{quote}
				
				Each of you bow your head and say, ``Yes.''

				That will do.

				Brethren, please arise.

				Each of you bring your right arm to the square.

				\begin{quote}
					You and each of you covenant and promise before God, angels, and these witnesses at this altar that you will observe and keep the law of chastity, as it has been explained to you.
				\end{quote}
				
				Each of you bow your head and say, ``Yes.''

				That will do.
				
			% ---
			\paragraph{The first token of the Melchizedek priesthood}
			% ---
				\label{TEMPLE-ENDOWMENT-ENDOW-PRE-1990-FIRST-MELCHIZEDEK}
				
				PETER: We will now give unto you the first token of the Melchizedek priesthood, or Sign of the Nail, with its accompanying name, sign, and penalty.

				\textit{[Peter demonstrates the token.]}

				We desire all to receive it. All arise.

				\textit{[All initiates receive the token. Peter then reveals the name, sign, and penalty for this token.]}

				I will now explain the covenant and obligation of secrecy which are associated with this token, its name, sign, and penalty, which you will be required to take upon yourselves.

				If I were receiving the endowment today, either for myself or for the dead, I would repeat in my mind these words, after making the sign, at the same time representing the execution of the penalty:

				I covenant in the name \rule{1cm}{0.15mm} that I will never reveal the first token of the Melchizedek priesthood, or Sign of the Nail, with its accompanying name, sign, and penalty. Rather than do so, I would suffer my life to be taken.
				
				\textit{[Peter leads the initiates in repeating this oath.]}

				That will do.

				We will return and report.

				\textit{[Peter, James, and John exit.]}
				
				\textit{[Elohim, Jehovah, and their messengers are heard speaking outside the Terrestrial Room.]}

				PETER: Jehovah, we have been down to the man Adam and his posterity, have placed the robe on the right shoulder, and have introduced them into the terrestrial world. We have put them under covenant to observe and keep the law of chastity. We have also given them the first token of the Melchizedek priesthood, or Sign of the Nail, with its accompanying name, sign, and penalty. This is our report.

				JEHOVAH: It is well, Peter, James, and John.

				Elohim---Peter, James, and John have been down to the man Adam and his posterity, have placed the robe on the right shoulder, and have introduced them into the terrestrial world. They have also put them under covenant to observe and keep the law of chastity. They have given unto them the first token of the Melchizedek priesthood, or Sign of the Nail, with its accompanying name, sign, and penalty. This is their report.

				ELOHIM: It is well.

				Jehovah, send down Peter, James, and John, and instruct them to give to the man Adam and his posterity in the terrestrial world the law of consecration, in connection with the law of the gospel and the law of sacrifice, and to cause them to receive it by covenant; to give unto them the second token of the Melchizedek priesthood, the patriarchal grip, or Sure Sign of the Nail, with its accompanying sign; and to teach them the order of prayer and prepare them in all things to receive further instructions at the veil. Then have them report at the veil.

				JEHOVAH: It shall be done, Elohim.

				Peter, James, and John, go down to the man Adam and his posterity in the terrestrial world and give unto them the law of consecration, in connection with the law of the gospel and the law of sacrifice, and cause them to receive it by covenant. Give unto them the second token of the Melchizedek priesthood, the patriarchal grip, or Sure Sign of the Nail, with its accompanying sign. Teach them the order of prayer and prepare them in all things to receive further instructions at the veil. Then report at the veil.

				PETER: It shall be done, Jehovah. Come, James and John, let us go down.
				
			% ---
			\paragraph{The law of consecration}
			% ---
				\label{TEMPLE-ENDOWMENT-ENDOW-PRE-1990-CONSECRATION}
				
				\textit{[Peter, James, and John return to the Terrestrial Room.]}

				PETER: A couple will now come to the altar.

				We are instructed to give unto you the law of consecration as contained in the book of Doctrine and Covenants, in connection with the law of the gospel and the law of sacrifice, which you have already received. It is that you do consecrate yourselves, your time, talents, and everything with which the Lord has blessed you, or with which he may bless you, to the Church of Jesus Christ of Latter-day Saints, for the building up of the kingdom of God on the earth and for the establishment of Zion.

				All arise.

				Each of you bring your right arm to the square.

				\begin{quote}
					You and each of you covenant and promise before God, angels, and these witnesses at this altar, that you do accept the law of consecration as contained in this, the book of Doctrine and Covenants \textit{[he displays the book]}, in that you do consecrate yourselves, your time, talents, and everything with which the Lord has blessed you, or with which he may bless you, to the Church of Jesus Christ of Latter-day Saints, for the building up of the kingdom of God on the earth and for the establishment of Zion.
				\end{quote}
				
				Each of you bow your head and say, ``Yes.''

				That will do.
				
			% ---
			\paragraph{The second token of the Melchizedek priesthood}
			% ---
				\label{TEMPLE-ENDOWMENT-ENDOW-PRE-1990-SECOND-MELCHIZEDEK}
				
				PETER: We will now give unto you the second token of the Melchizedek priesthood, the patriarchal grip, or Sure Sign of the Nail, with its accompanying sign.

				\textit{[Peter demonstrates the token.]}

				We desire all to receive it. All arise.

				\textit{[All initiates receive the token.]}

				This token has a name and a sign, but no penalty is mentioned. However, you will be under just as sacred an obligation of secrecy in connection with this token and sign as you are with the other tokens and signs of the holy priesthood which you have received in the temple this day.

				The name of this token will not be given unto you at this stage in the endowment, but it will be given later on.

				\textit{[Peter demonstrates the sign for this token and then leads initiates in making the sign.]}

				That will do.
				
			% ---
			\paragraph{The order of prayer}
			% ---
				\label{TEMPLE-ENDOWMENT-ENDOW-PRE-1990-ORDER-PRAYER}
				
				PETER: Brethren and sisters, with the robe on the right shoulder you are prepared to be taught the true order of prayer and to be introduced at the veil.

				A few couples will please come forward and form a circle around the altar.

				\textit{[The couples are arranged so that men and women alternate around the circle. James directs the prayer circle.]}

				JAMES: Only the best of feelings should exist in the circle. If any of you have unkind feelings toward any member of this circle, you are invited to withdraw so that the Spirit of the Lord may be unrestrained.

				In the circle we make the signs of all the tokens of the holy priesthood.

				\textit{[James leads those in the circle in making each sign, in succession. As they do this, he reminds initiates of the name for each token, except the second token of the Melchizedek priesthood, the name of which has not yet been revealed.]}

				We have here a list of names of persons who are sick or otherwise afflicted, whom we are requested to remember in our prayer. We will place the list upon the altar and request the faith of those present in behalf of these persons.

				The sisters in the room will please veil their faces.

				Each brother in the circle will take the sister at his left by the right hand in the patriarchal grip.

				Each of you bring your left arm to the square and rest it upon the shoulder or arm of the person at your left.

				The brethren and sisters in the circle will repeat the words of the prayer.

				\textit{[James kneels before the altar, within the circle, makes the sign of the second token of the Aaronic priesthood, and offers an extemporaneous prayer. He pauses after each sentence, so that those in the circle may repeat his words.]}

				PETER: The sisters will unveil their faces, and the brethren and sisters in the circle will return to their seats.

				We will now uncover the veil.
				
		% -----
		\subsubsection{The Veil}
		% -----
			\label{TEMPLE-ENDOWMENT-ENDOW-PRE-1990-VEIL}
			
			PETER: Brethren and sisters, I will now explain the marks on the veil.

			These four marks are the marks of the holy priesthood, and corresponding marks are found in your individual garment.

			This one on the right is the mark of the square. It is placed in the garment over the right breast, suggesting to the mind exactness and honor in keeping the covenants entered into this day.

			This one on the left is the mark of the compass. It is placed in the garment over the left breast, suggesting to the mind an undeviating course leading to eternal life; a constant reminder that desires, appetites, and passions are to be kept within the bounds the Lord has set; and that all truth may be \hyperref[CIRCUMSCRIBE]{circumscribed} into one great whole.

			This is the navel mark. It is placed in the garment over the navel, suggesting to the mind the need of constant nourishment to body and spirit.

			This is the knee mark. It is placed in the right leg of the garment so as to be over the kneecap, suggesting that every knee shall bow and every tongue shall confess that Jesus is the Christ.

			These other three marks are for convenience in working at the veil. Through this one, the person representing the Lord%
				\footnote{%
					\label{ENDOWMENT-lord}%
					So the person on the other side of the veil is supposed to represent the ``Lord,'' or Christ, I think. See \hyperref[2NEPHI-9-41]{2 Nephi 9:41}. This seems to represent some kind of personal interaction we have with Christ across the veil---as Peter reports to Jehovah later, ``They are now ready to converse with the Lord through the veil.'' Is this before the resurrection? After? I'm not sure.
					}
			puts forth his right hand to test our knowledge of the tokens of the holy priesthood. Through the one on our right, he asks us certain questions; through the one on the left, we give our answers.
			
			\textit{[Peter stands before the veil, in view of the temple initiates. A person representing the Lord stands concealed behind the veil.]}

			PETER: As all of you will have to pass through the veil, we will show you how this is to be done.

			\textit{[Peter stands before a part in the veil, off to the side of the marks whose meaning he has just explained.]}

			PETER: The person is brought to this point, and the worker gives three distinct taps with the mallet...

			\textit{[Peter does so. The person representing the Lord puts his hand through the part in the veil, opening it slightly.]}

			...whereupon the Lord parts the veil, and asks:

			LORD: What is wanted?

			PETER: Adam, having been true and faithful in all things, desires further light and knowledge by conversing with the Lord through the veil.

			LORD: Present him at the veil, and his request shall be granted.

			\textit{[Peter stands in front of the marks on the veil. The person representing the Lord puts his hand through the mark which Peter has indicated for this purpose.]}

			PETER: The person is then brought to this point, whereupon the Lord puts forth his right hand, gives the first token of the Aaronic priesthood, and asks:

			LORD: What is that?

			PETER: The first token of the Aaronic priesthood.

			LORD: Has it a name?

			PETER: It has.

			LORD: Will you give it to me?

			PETER: I will, through the veil.

			The person then gives, through the veil, the name of this token, which is the new name received in the temple today.

			PETER: The Lord then gives the second token of the Aaronic priesthood and asks:

			LORD: What is that?

			PETER: The second token of the Aaronic priesthood.

			LORD: Has it a name?

			PETER: It has.

			LORD: Will you give it to me?

			PETER: I will, through the veil.

			The person then gives the name of this token, which is \rule{1cm}{0.15mm}.

			PETER: The Lord then gives the first token of the Melchizedek priesthood, or Sign of the Nail, and asks:

			LORD: What is that?

			PETER: The first token of the Melchizedek priesthood, or Sign of the Nail.

			LORD: Has it a name?

			PETER: It has.

			LORD: Will you give it to me?

			PETER: I will, through the veil.

			The person then gives the name of this token, which is \rule{1cm}{0.15mm}.

			PETER: The Lord then gives the second token of the Melchizedek priesthood, the patriarchal grip, or Sure Sign of the Nail, and asks:

			LORD: What is that?

			PETER: The second token of the Melchizedek priesthood, the patriarchal grip, or Sure Sign of the Nail.

			LORD: Has it a name?

			PETER: It has.

			LORD: Will you give it to me?

			PETER: I cannot. I have not yet received it. For this purpose I have come to converse with the Lord through the veil.

			LORD: You shall receive it, upon the five points of fellowship, through the veil.

			PETER: The five points of fellowship are: inside of right foot by the side of right foot, knee to knee, breast to breast, hand to back, and mouth to ear.

			The Lord then gives the name of this token and asks:

			LORD: What is that?

			PETER: The second token of the Melchizedek priesthood, the patriarchal grip, or Sure Sign of the Nail.

			LORD: Has it a name?

			PETER: It has.

			LORD: Will you give it to me?

			PETER: I will, upon the five points of fellowship, through the veil.

			The person then repeats back to the Lord the name of this token as he received it, whereupon the Lord says:

			LORD: That is correct.

			\textit{[The person representing the Lord withdraws his hand from view.]}

			\textit{[Peter stands again before the part in the veil.]}

			PETER: The person is again brought to this point, and the worker gives three distinct taps with the mallet.

			\textit{[Peter does so.]}

			The Lord parts the veil and asks:

			LORD: What is wanted?

			PETER: Adam, having conversed with the Lord through the veil, desires now to enter his presence.

			The Lord puts forth his right hand, takes the person by the right hand, and says:

			LORD: Let him enter.

			PETER: He is admitted into the presence of the Lord.

			\textit{[The person representing the Lord pulls Peter partway through the veil to demonstrate.]}

			We will now report.

			\textit{[Peter, James, and John pass behind the veil.]}
			
			\textit{[Elohim, Jehovah, and their messengers are heard speaking behind the veil.]}

			PETER: Jehovah, we have been down to the man Adam and his posterity in the terrestrial world, and have given unto them the law of consecration, and have caused them to receive it by covenant. We have given unto them the second token of the Melchizedek priesthood, the patriarchal grip, or Sure Sign of the Nail, with its accompanying sign, and have taught them the order of prayer. They are now ready to converse with the Lord through the veil. This is our report.

			JEHOVAH: It is well, Peter, James, and John.

			Elohim--Peter, James, and John have been down to the man Adam and his posterity in the terrestrial world and have done all that they were commanded to do.

			ELOHIM: It is well.

			Jehovah, instruct Peter, James, and John to introduce the man Adam and his posterity at the veil, where we will give unto them the name of the second token of the Melchizedek priesthood, the patriarchal grip, or Sure Sign of the Nail, preparatory to their entering into our presence.

			JEHOVAH: It shall be done, Elohim.

			Peter, James, and John, you will introduce Adam and his posterity at the veil, where we will give unto them the name of the second token of the Melchizedek priesthood, the patriarchal grip, or Sure Sign of the Nail, preparatory to their entering into our presence.

			PETER: It shall be done, Jehovah. Come, James and John, we will introduce them at the veil.

			\textit{[Peter, James, and John return from behind the veil.]}

			Brethren and sisters, we are instructed to introduce you at the veil, where you will receive the name of the second token of the Melchizedek priesthood, the patriarchal grip, or Sure Sign of the Nail, preparatory to your entering into the presence of the Lord.
			
			% --
			\paragraph{The lecture at the veil} % *LECTURE
			% --
				\label{TEMPLE-ENDOWMENT-ENDOW-PRE-1990-LECTURE}
				
				PETER: A lecture%
					\footnote{%
						\label{ENDOWMENT-lecture}%
						Remember that the original ``lecture at the veil'' was given by BY for the St. George temple, and involved Adam-God. See the typescript version of L. John Nuttall's diary, for 7 February 1877, starting on apge 18 (PDF 137), or my transcription \hyperref[EMS-1877-BY-7-FEB-1877]{here}.
						}
				will now be given which summarizes the instructions, ordinances, and covenants, and also the tokens, with their keywords, signs, and penalties, pertaining to the endowment, which you have thus far received.

				You should try to remember and keep in mind all that you have heard and seen, and may yet hear and see, in this house. The purpose of this lecture is to assist you to remember that which has been taught you this day. You must keep in mind that you are under a solemn obligation never to speak outside of the temple of the Lord of the things you see and hear in this sacred place.

				LECTURER: Brethren and sisters, the ordinances of the endowment as here administered, long withheld from the children of men, pertain to the dispensation of the fullness of time and have been revealed to prepare the people for exaltation in the celestial kingdom, where God and Christ dwell.

				The deep meaning of the eternal truths constituting the endowment has been set forth in brief instructions and by symbolic representation. If you give prayerful and earnest thought to the holy endowment, you will obtain the understanding and spirit of the work done in the temples of the Lord. The privilege of laboring here for the dead permits us to enter the temple frequently, and to refresh our memories, and to enlarge our understanding of the endowment.

				You were first washed and anointed, a garment was placed upon you, and a new name was given you. This name you should always remember; but you must never reveal it to any person, except at the veil.

				You then entered [the Creation Room]. [There] you heard the voices of persons representing a council of the Gods---Elohim, Jehovah, and Michael. Elohim said: ``See, yonder is matter unorganized. Go ye down and organize it into a world, like unto the other worlds that we have heretofore formed.'' As the creation of the earth progressed, you heard the commands and the reports of the persons representing the Gods.

				If we are faithful, we shall enter the celestial kingdom and there hear and know the Gods of heaven. They are perfect; we are imperfect. They are exalted; we may attain exaltation.

				Our spirits at one time lived with the Gods; but each of us was given the privilege of coming upon this earth to take upon himself a body, so that the spirit might have a house in which to dwell.

				Michael, one of the council of the Gods, became the man Adam, to whom was given the woman Eve. However, as Adam he did not remember his life and labors in the council. It is so with us all. We came into the world with no memory of our previous existence.

				We then followed Adam and Eve into the garden, where Elohim provided that they might eat freely of all kinds of fruit of the garden, except the fruit of the tree of knowledge of good and evil. He forbade them to partake of this fruit, saying that in the day they did so, they should surely die.

				When Adam and Eve were left alone in the garden, Satan appeared and tempted them. Eve yielded to the temptation, partook of the fruit, and offered it to Adam. Adam had resisted the temptation of Satan, but when Eve offered him the forbidden fruit, he partook of it that they might continue together and perpetuate the human race.

				Adam and Eve now understood that it was Lucifer who had tempted them. They became self-conscious. Discovering their nakedness and hearing the voice of the Lord, they made aprons of fig leaves and hid themselves. They had learned that everything has its opposite, such as good and evil, light and darkness, pleasure and pain.

				The Lord again entered the garden. Adam and Eve confessed their disobedience. The Lord cursed Satan and cast him out of the garden of Eden, and the Lord commanded: ``Let cherubim, and a flaming sword be placed to guard the way of the tree of life, lest Adam put forth his hand, and partake of the fruit thereof, and live forever in his sins.''

				Before their departure, however, instructions were given them. Addressing Eve, the Lord said: ``Because thou hast hearkened to the voice of Satan and hast partaken of the forbidden fruit and given unto Adam, I will greatly multiply thy sorrow and thy conception. In sorrow shalt thou bring forth children; nevertheless, thou mayest be preserved in childbearing. Thy desire shall be to thy husband, and he shall rule over thee in righteousness.''

				To Adam, the Lord said: ``Because thou hast hearkened to the voice of thy wife and hast partaken of the forbidden fruit, the earth shall be cursed for thy sake. Instead of producing fruits and flowers spontaneously, it shall bring forth thorns, thistles, briars, and noxious weeds to afflict and torment man; and by the sweat of thy face shalt thou eat thy bread all the days of thy life, for dust thou art, and unto dust shalt thou return.''

				Having been commanded, Jehovah provided Adam and Eve with coats of skins for a covering. The garment which was placed upon you after you had been washed and anointed represents the coat of skins, or covering, of Adam and Eve. They were also promised that further light and knowledge would be given them.

				The law of obedience was then taught Adam and Eve and accepted by them. Eve covenanted with Adam that thenceforth she would obey the law of her husband and abide by his counsel in righteousness; and Adam covenanted with the Lord that he would obey the Lord and keep his commandments. You, likewise, covenanted to comply with the law of obedience.

				The law of sacrifice accompanying the law of obedience, as contained in the Old and New Testaments of the Bible, was next presented to Adam, and you were all placed under covenant to observe it.

				The law of obedience and sacrifice includes the promise of the Savior, the Only Begotten of the Father, who is full of grace and truth and who by his sacrifice has become the Redeemer of mankind. All things should be done in the name of the Son. An angel of the Lord explained this to Adam, who was given the privilege of showing his obedience by offering sacrifices to the Lord in similitude of the sacrifice of Jesus Christ. Later, the people of Israel lived under this law, which continued in force until the death of Jesus Christ.

				The first token of the Aaronic priesthood, with its accompanying name, sign, and penalty, was given you; and you were told that the name of this token is your new name, or the new name of the dead, if officiating for the dead. The sacred nature of the tokens of the priesthood was carefully explained at this time. You were placed under solemn covenant never to reveal these tokens, with their accompanying names, signs, and penalties, even at the peril of your life. You were told that the execution of the penalties indicate different ways in which life may be taken.

				Then Adam and Eve were driven out of the garden into the telestial kingdom, or the lone and dreary world, the world in which we are now living. There Adam offered a prayer saying, ``Oh God, hear the words of my mouth,'' repeating it three times.

				Satan entered and, claiming to be the god of this world, asked Adam what he desired. Adam replied that he was waiting for messengers from his Father. Satan declared that a preacher would soon arrive. A man representing a sectarian minister entered and preached doctrine, which Adam did not accept.

				Peter, James, and John were sent down by the Lord to learn, without disclosing their identity, if the man Adam had been faithful to his covenants. They found that he had been faithful and so reported.

				They were sent down again, this time in their true character as apostles of the Lord Jesus Christ, to visit and to instruct Adam and his posterity in the telestial world. Before so teaching the people, they cast Satan out.

				The law of the gospel, as contained in the Book of Mormon and the Bible, was then given Adam and his posterity. You were placed under covenant to obey the law of the gospel, and to avoid all lightmindedness, loud laughter, evil speaking of the Lord's anointed, and taking the name of the Lord in vain.

				The robe of the holy priesthood was placed upon your left shoulder, according to the order of the Aaronic priesthood.

				The second token of the Aaronic priesthood was given you, with its name, sign, and penalty, and you were informed that the name of this token is \rule{1cm}{0.15mm}.

				The robe of the holy priesthood was then changed to the right shoulder, as was done anciently when officiating in the ordinances of the Melchizedek priesthood. With the robe on the right shoulder, you have authority also, if called to the bishopric, to act in the Aaronic priesthood.

				You were then introduced, with the robe of the holy priesthood on the right shoulder, into the terrestrial kingdom. The law of chastity was there explained to you in plainness, and you were placed under covenant to obey this law.

				The first token of the Melchizedek priesthood, or Sign of the Nail, with its accompanying name, sign, and penalty, was next given you. You were told that the name of the first token of the Melchizedek priesthood is \rule{1cm}{0.15mm}.

				The book of Doctrine and Covenants, in connection with the Book of Mormon and the Bible, was presented to you; and the law of consecration, as contained in the book of Doctrine and Covenants, was explained to you; and you received this law by covenant.

				The second token of the Melchizedek priesthood, the patriarchal grip, or Sure Sign of the Nail, or the Nail in the Sure Place,%
					\footnote{%
						% \label{}%
						See \hyperref[ISAIAH-22-20]{Isaiah 22:20--25}.
						}
				was given you, together with its sign. The name of this token will be given you at the veil.

				This token has reference to the crucifixion of the Savior. When he was placed upon the cross, the crucifiers drove nails through the palms of his hands; then, fearing that the weight of his body would cause the nails to tear through the flesh of his hands, they drove nails through his wrists. Hence in the palm is the Sign of the Nail, and in the wrist is the Sure Sign of the Nail, or the Nail in the Sure Place.

				You have now progressed so far in the endowment that you are ready to receive the name of the second token of the Melchizedek priesthood and to pass through the veil into the celestial kingdom.

				The sisters in this company who are to be married and sealed for time and eternity should be taken through the veil by their intended husband. Others will be taken through the veil by the regular temple workers.

				Brethren and sisters, you will have received this day the sacred ordinances of the endowment. The eternal plan of salvation for man, as he journeys from his pre-existent state to his future high place in the celestial kingdom, has been presented to you. You have covenanted to obey all the laws of the gospel, including the laws of obedience, sacrifice, chastity, and consecration, which make possible an exaltation with the Gods; and you have received the first and second tokens of the Aaronic priesthood and the first and second tokens of the Melchizedek priesthood, with the names, signs, and penalties of these tokens, except the name of the second token of the Melchizedek priesthood, which will be given you at the veil.

				All this is done for the glory, honor, and endowment of the children of Zion.

				Brethren and sisters, strive to comprehend the glorious things presented to you this day. No other people on earth have ever had this privilege, except as they have received the keys of the priesthood given in the endowment.

				These are what are termed the mysteries of godliness---that which will enable you to understand the expression of the Savior, made just prior to his betrayal: ``This is life eternal, that they might know thee, the only true God, and Jesus Christ, whom thou has sent.''

				May God bless you all. Amen.
				
			% --
			\paragraph{Passing through the veil}
			% --
				\label{TEMPLE-ENDOWMENT-ENDOW-PRE-1990-PASSING}
				
				\textit{[Initiates are presented individually at the veil. The veil worker taps three times with the mallet.]}

				LORD: What is wanted?

				WORKER: Adam, having been true and faithful in all things, desires further light and knowledge by conversing with the Lord through the veil \textit{[for and in behalf of \rule{1cm}{0.15mm}, who is dead]}.

				LORD: Present him at the veil, and his request shall be granted.

				LORD: What is that?

				INITIATE: The first token of the Aaronic priesthood.

				LORD: Has it a name?

				INITIATE: It has.

				LORD: Will you give it to me?

				INITIATE: I will, through the veil.

				\textit{[The initiate gives the token's name.]}

				LORD: What is that?

				INITIATE: The second token of the Aaronic priesthood.

				LORD: Has it a name?

				INITIATE: It has.

				LORD: Will you give it to me?

				INITIATE: I will, through the veil.

				\textit{[The initiate gives the token's name.]}

				LORD: What is that?

				INITIATE: The first token of the Melchizedek priesthood, or Sign of the Nail.

				LORD: Has it a name?

				INITIATE: It has.

				LORD: Will you give it to me?

				INITIATE: I will, through the veil.

				\textit{[The initiate gives the token's name.]}

				LORD: What is that?

				INITIATE: The second token of the Melchizedek priesthood, the patriarchal grip, or Sure Sign of the Nail.

				LORD: Has it a name?

				INITIATE: It has.

				LORD: Will you give it to me?

				INITIATE: I cannot. I have not yet received it. For this purpose I have come to converse with the Lord through the veil.

				LORD: You shall receive it, upon the five points of fellowship, through the veil.

				\textit{[Still holding the patriarchal grip, the initiate and the person representing the Lord embrace upon the five points of fellowship through the veil. The initiate's left arm passes through the mark of the compass, and the left arm of the person representing the Lord passes through the mark of the square. The person representing the Lord speaks into the initiate's ear the name of the token.]}%
					\footnote{%
						% \label{}%
						``Health in the navel, marrow in the bones, strength in the loins and in the sinews. Power in the priesthood be upon me and upon my posterity through all generations of time, and throughout all eternity.'' See \hyperref[PROVERBS-3-8]{Proverbs 3:8}; \hyperref[DC89-18]{DC 89:18}; \hyperref[DC109-24]{DC 109:24}; 
						}

				LORD: What is that?

				INITIATE: The second token of the Melchizedek priesthood, the patriarchal grip, or Sure Sign of the Nail.

				LORD: Has it a name?

				INITIATE: It has.

				LORD: Will you give it to me?

				INITIATE: I will, upon the five points of fellowship, through the veil.

				\textit{[The initiate speaks back the name of the token.]}

				LORD: That is correct.

				\textit{[The initiate returns to the part in the veil; the worker taps three times with the mallet.]}

				LORD: What is wanted?

				WORKER: Adam, having conversed with the Lord through the veil, desires now to enter his presence.

				LORD: Let him enter.

				\textit{[The initiate is brought through the veil into the Celestial Room.]}